\chapter{Il caso discreto: da finito a numerabile}

Fino ad ora abbiamo considerato solo spazi di probabilità finiti. Questo è un limite significativo per descrivere molti fenomeni naturali. Ad esempio, non potremmo descrivere una variabile aleatoria che conta il numero di lanci necessari per ottenere testa, poiché i lanci potrebbero essere teoricamente infiniti, richiedendo un codominio come $\N$. Allo stesso modo, non potremmo descrivere processi su orizzonti temporali infiniti. 

Dobbiamo estendere il nostro quadro teorico dai casi finiti ai casi infiniti, in particolare spazi numerabili (o più grandi). Ecco le estensioni chiave che affronteremo:
\begin{itemize}
    \item \textbf{Spazio campionario} $\Omega$: da finito a infinito.
    \item \textbf{$\sigma$-algebra}: chiusura rispetto alle unioni da finita a \emph{numerabile}.
    \item \textbf{Misura di probabilità}: dall'additività finita all'\emph{additività numerabile} ($\sigma$-additività).
    \item \textbf{Famiglie di variabili aleatorie}: da collezioni finite a collezioni \emph{numerabili} $\{X_j\}_{j \in \N}$.
\end{itemize}

\section{Spazi di probabilità generali}

\begin{definition}[$\sigma$-algebra]
Una collezione $\F \subseteq 2^\Omega$ è chiamata \emph{$\sigma$-algebra} su $\Omega$ se:
\begin{enumerate}
    \item $\Omega \in \F$,
    \item $A \in \F \implies A^c \in \F$ (chiusura rispetto al complemento),
    \item Per ogni collezione numerabile $\{A_n\}_{n \in \N} \subseteq \F$, l'unione $\bigcup_{n \in \N} A_n \in \F$ (chiusura rispetto all'unione numerabile).
\end{enumerate}
La coppia $(\Omega, \F)$ è detta \emph{spazio misurabile}.
\end{definition}

\begin{definition}[Misura di probabilità ($\sigma$-additiva)]
Sia $\F$ una $\sigma$-algebra. Una funzione $\PP : \F \to [0, 1]$ è una \emph{misura di probabilità} su $(\Omega, \F)$ se:
\begin{enumerate}
    \item $0 \le \PP(A) \le 1$ per ogni $A \in \F$,
    \item $\PP(\Omega) = 1$,
    \item \textbf{$\sigma$-additività}: per ogni collezione numerabile di eventi \emph{disgiunti} a due a due $\{A_n\}_{n \in \N}$, vale:
    \[ \PP\left(\bigcup_{n=1}^\infty A_n\right) = \sum_{n=1}^\infty \PP(A_n) \]
\end{enumerate}
\end{definition}

\begin{spiegazione}
    Il passaggio dall'additività finita alla $\sigma$-additività (numerabile) è il cuore della probabilità moderna introdotta da Kolmogorov. Ci permette di sommare probabilità di infiniti eventi disgiunti e garantisce che la probabilità si comporti bene quando facciamo limiti di eventi (continuità dal basso e dall'alto).
\end{spiegazione}

\section{Variabili aleatorie discrete}

Limitiamo la nostra teoria alle variabili aleatorie che hanno un codominio \emph{al più numerabile}. Queste sono chiamate \textbf{variabili aleatorie discrete}. Tutti i concetti di densità discreta e distribuzione del Capitolo 1 si estendono naturalmente sostituendo le somme finite con le \emph{serie} numerabili.

\subsection{Distribuzione di Poisson}

Una distribuzione fondamentale sugli interi naturali $\N$ è la distribuzione di Poisson.

\begin{definition}[Distribuzione di Poisson]
La distribuzione di Poisson di parametro $\lambda > 0$, indicata con $\mathcal{P}(\lambda)$, è la legge di probabilità su $\N$ definita da:
\[ \PP(X = k) = e^{-\lambda} \frac{\lambda^k}{k!}, \quad k \in \N \]
\end{definition}

\begin{hint}[Somma di Poisson]
    Ricorda che se sommiamo due variabili di Poisson indipendenti con parametri $\lambda_1$ e $\lambda_2$, il risultato è ancora una variabile di Poisson con parametro $\lambda_1 + \lambda_2$.
\end{hint}

\section{Schema di Bernoulli e sequenze indipendenti}

Vogliamo modellare un esperimento che consiste in infiniti lanci indipendenti di una moneta. Abbiamo bisogno della garanzia matematica che esista uno spazio di probabilità abbastanza grande da supportare una sequenza infinita di variabili indipendenti.

\begin{theorem}[Teorema di estensione di Kolmogorov - Schema di Bernoulli]
Sia $(E, \mathcal{E})$ uno spazio discreto e $\mu$ una distribuzione di probabilità. Esiste uno spazio di probabilità $(\Omega, \F, \PP)$ e una sequenza \textbf{infinita} di variabili aleatorie \textbf{indipendenti e identicamente distribuite (i.i.d.)} $X_1, X_2, \dots$ a valori in $E$ tale che ogni $X_i$ abbia distribuzione $\mu$.
\end{theorem}

Questo teorema legittima l'uso di infinite prove ripetute (lo schema di Bernoulli).

\subsection{Distribuzione Geometrica}

Consideriamo una sequenza infinita di tentativi di Bernoulli i.i.d.\ con probabilità di successo $p \in (0,1)$. Sia $T$ la variabile aleatoria che conta il \emph{numero di tentativi fino al primo successo} (incluso). 

\begin{definition}[Distribuzione Geometrica]
La legge del "tempo di primo successo" $T$ si chiama \emph{distribuzione geometrica di parametro $p$}. La sua densità discreta è:
\[ \PP(T = n) = (1-p)^{n-1} p, \quad n \in \N, n \ge 1 \]
\end{definition}
(Notiamo che la probabilità che non ci sia mai un successo è zero, $\PP(T=\infty)=0$).

\begin{theorem}[Proprietà di Assenza di Memoria]
La distribuzione geometrica è l'unica distribuzione discreta ad avere la proprietà di assenza di memoria (\emph{memoryless property}):
\[ \PP(T = n + k | T > n) = \PP(T = k) \quad \forall n, k \ge 1 \]
\end{theorem}

\begin{spiegazione}
    L'assenza di memoria significa che "la moneta non ha memoria". Se hai appena fatto 10 lanci fallimentari e aspetti il successo, la probabilità di dover aspettare altri 3 lanci è IDENTICA alla probabilità che avresti avuto all'inizio di dover aspettare 3 lanci. Il passato non cambia l'attesa futura!
\end{spiegazione}

\section{Valore Atteso nel Caso Discreto Generale}

Quando passiamo a spazi numerabili, le somme diventano serie infinite. Dobbiamo fare attenzione alla convergenza. 

\begin{definition}[Variabile aleatoria semi-integrabile e valore atteso]
Sia $X : \Omega \to \R$ una variabile aleatoria discreta (a valori in un insieme numerabile). Definiamo il suo \emph{valore atteso} separando la parte positiva $X^+$ dalla parte negativa $X^-$:
\[ \E[X] = \E[X^+] - \E[X^-] \]
Se almeno uno dei due termini è finito, diciamo che $X$ è \emph{semi-integrabile}. Se entrambi sono finiti, $X$ è \emph{integrabile} e scriviamo $X \in L^1(\Omega, \F, \PP)$.
\end{definition}

Le proprietà fondamentali del valore atteso (Linearità, Monotonia, Positività, Legge dello Statistico Inconsapevole) rimangono identiche, a patto di garantire la convergenza assoluta delle serie coinvolte (cioè $X, Y \in L^1$).

Un risultato estremamente utile per variabili a valori interi positivi è il seguente:
\begin{theorem}[Formula alternativa per il valore atteso intero]
Se $X$ assume valori in $\N \cup \{0, \infty\}$, il suo valore atteso può essere calcolato usando le probabilità cumulate complementari (code della distribuzione):
\[ \E[X] = \sum_{n=1}^\infty \PP(X \ge n) \]
\end{theorem}

\subsection{Varianza e Disuguaglianze Notevoli}
La Varianza si definisce come $\operatorname{Var}(X) = \E[(X-\E[X])^2]$ per variabili con varianza finita. In questo contesto, le disuguaglianze di Markov e Chebyshev sono strumenti fondamentali per stimare probabilità anche senza conoscere la densità esatta.

\begin{theorem}[Disuguaglianza di Markov e di Chebyshev]
\textbf{Markov:} Per ogni $X$ discreta non negativa e per ogni $a > 0$:
\[ \PP(X \ge a) \le \frac{\E[X]}{a} \]

\textbf{Chebyshev:} Per ogni $X$ discreta con media $m$ e varianza $\sigma^2$ finita, e per ogni $\epsilon > 0$:
\[ \PP(|X - m| \ge \epsilon) \le \frac{\operatorname{Var}(X)}{\epsilon^2} \]
\end{theorem}

\section{Valore Atteso Condizionato (Caso Generale)}

Anche il valore atteso condizionato $\E[X|Y]$ (inteso come variabile aleatoria) mantiene le sue proprietà fondamentali, con l'accortezza di sostituire le somme finite con le serie.
Ricordiamo le proprietà più importanti, dove si assume che tutte le variabili siano in $L^1$:
\begin{enumerate}
    \item \textbf{Legge delle aspettative iterate (Tower Property):} $\E[\E[X|Y]] = \E[X]$.
    \item \textbf{Tirar fuori ciò che è noto (Taking out what is known):} $\E[X Z | Y] = Z \E[X | Y]$ se $Z$ è una funzione deterministica di $Y$.
    \item \textbf{Disuguaglianza di Jensen Condizionata:} Se $\phi$ è convessa e $X \in L^1$, allora $\phi(\E[X|Y]) \le \E[\phi(X)|Y]$.
    \item \textbf{Freezing Lemma:} $\E[\phi(X,Y)|Y](\omega) = \E[\phi(X, y)]\big|_{y=Y(\omega)}$.
\end{enumerate}

\begin{hint}[La formula di disintegrazione e il Freezing Lemma]
    Combinando il Freezing Lemma con la Tower Property otteniamo una generalizzazione della formula di disintegrazione:
    \[ \E[\phi(X,Y)] = \E[\E[\phi(X,Y)|Y]] \]
    Se $X$ e $Y$ sono indipendenti, calcolare la media interna è facilissimo perché $\E[\phi(X,y)|Y=y] = \E[\phi(X,y)]$.
\end{hint}

\begin{esempio}[Equazione di Wald]
Sia $(X_i)_{i \ge 1}$ una successione di variabili i.i.d. con media $\mu$, e sia $N$ una variabile a valori in $\N$, \emph{indipendente} dalle $X_i$. Definiamo la somma aleatoria $S_N = \sum_{i=1}^N X_i$.
Qual è il suo valore atteso? Usando il condizionamento rispetto a $N$:
\[ \E[S_N|N=n] = \E\left[\sum_{i=1}^n X_i\right] = n\mu \]
Quindi $\E[S_N|N] = N\mu$. Applicando la Tower Property: $\E[S_N] = \E[N\mu] = \mu \E[N]$. Questa è nota come Equazione di Wald.
\end{esempio}

\section{Convergenze di Variabili Aleatorie}

Una delle necessità più grandi nel passare a collezioni numerabili di variabili (es. le medie di $n$ esperimenti per $n \to \infty$) è poter studiare il limite delle variabili aleatorie.
In probabilità, abbiamo tre concetti principali di convergenza.

\begin{definition}[Tipi di Convergenza]
Data una successione $X_n$ e una variabile limite $X$:
\begin{enumerate}
    \item \textbf{Convergenza Quasi Certamente (a.s.):} $X_n \to X$ q.c. se $\PP(\{\omega \in \Omega : \lim_{n\to\infty} X_n(\omega) = X(\omega)\}) = 1$. (Per quasi ogni esito, la sequenza numerica converge).
    \item \textbf{Convergenza in Probabilità:} $X_n \stackrel{P}{\to} X$ se per ogni $\epsilon > 0$, $\lim_{n\to\infty} \PP(|X_n - X| \ge \epsilon) = 0$. (La probabilità che $X_n$ disti da $X$ più di $\epsilon$ svanisce).
    \item \textbf{Convergenza in Legge (o in Distribuzione):} $X_n \stackrel{\mathcal{L}}{\to} X$ se per ogni funzione $\phi$ limitata e continua, $\E[\phi(X_n)] \to \E[\phi(X)]$. (Le \emph{distribuzioni} di $X_n$ diventano simili a quella di $X$, ma le variabili stesse su $\Omega$ potrebbero non essere "vicine").
\end{enumerate}
\end{definition}

\begin{spiegazione}
    Esiste una gerarchia rigida tra queste convergenze:
    \[ \text{Quasi Certamente} \implies \text{In Probabilità} \implies \text{In Legge} \]
    Le implicazioni inverse \textbf{non} valgono in generale. 
    Ad esempio, se lancio all'infinito due dadi sempre diversi ma entrambi bilanciati, le loro distribuzioni (leggi) sono identiche, ma il risultato del dado $n$-esimo non "converge numericamente" in probabilità al risultato di un dado $X$.
\end{spiegazione}

\subsection{Teoremi di convergenza in media}

Passare al limite dentro un valore atteso ($\lim \E[X_n] = \E[\lim X_n]$) non è sempre possibile. Esistono tre grandi teoremi per legittimare questo passaggio, mutuati dalla teoria dell'integrazione di Lebesgue:

\begin{theorem}[Teorema di Convergenza Monotona (Beppo-Levi)]
Se $X_n \ge 0$ e $X_n \uparrow X$ monotonamente (quasi certamente), allora:
\[ \lim_{n\to\infty} \E[X_n] = \E[X] \]
\end{theorem}

\begin{theorem}[Lemma di Fatou]
Se $X_n \ge 0$ (quasi certamente) per ogni $n$, allora l'aspettativa del limite inferiore è minore o uguale al limite inferiore delle aspettative:
\[ \E\left[\liminf_{n\to\infty} X_n\right] \le \liminf_{n\to\infty} \E[X_n] \]
\end{theorem}

\begin{theorem}[Teorema di Convergenza Dominata (Lebesgue)]
Se $X_n \to X$ (in probabilità o quasi certamente) e se esiste una variabile "dominante" $G \in L^1$ tale che $|X_n| \le G$ per ogni $n$, allora:
\[ \lim_{n\to\infty} \E[X_n] = \E[X] \]
\end{theorem}

\section{La Legge dei Grandi Numeri}

Uno dei risultati più celebri della teoria della probabilità è la \emph{Legge dei Grandi Numeri}. In senso lato, afferma che:
\emph{"La media empirica di una successione di variabili aleatorie identicamente distribuite converge (in un senso opportuno) al loro valore atteso comune."}

Esistono due versioni principali, a seconda del tipo di convergenza (e delle ipotesi richieste).

\begin{theorem}[Legge debole dei grandi numeri]
Sia $(X_n)_{n \in \N}$ una successione di variabili i.i.d. discrete a valori reali tali che $\E[X_1^2] < \infty$ (il che implica varianza finita $\sigma^2$). Sia $m = \E[X_1]$. Posto:
\[ \bar{S}_n = \frac{X_1 + \cdots + X_n}{n} \]
allora per ogni $\epsilon > 0$:
\[ \lim_{n\to\infty} \PP(|\bar{S}_n - m| \ge \epsilon) = 0 \]
(cioè $\bar{S}_n \stackrel{P}{\to} m$).
\end{theorem}
\emph{Dimostrazione (Idea):} Usando la linearità, $\E[\bar{S}_n] = m$. Per l'indipendenza, $\operatorname{Var}(\bar{S}_n) = \frac{\sigma^2}{n}$. Applicando Chebyshev, $\PP(|\bar{S}_n - m| \ge \epsilon) \le \frac{\sigma^2}{n\epsilon^2}$, che tende a 0.

\begin{theorem}[Legge forte dei grandi numeri di Borel]
Sia $(X_n)_{n \in \N}$ i.i.d. e \textbf{limitate} (esiste $C$ tale che $|X_i| \le C$). Posto $m = \E[X_1]$, si ha:
\[ \PP\left( \lim_{n\to\infty} \frac{S_n}{n} = m \right) = 1 \]
(cioè $\bar{S}_n \to m$ quasi certamente).
\end{theorem}

\begin{hint}[Riflessione Frequenzista e Assiomatica]
È fondamentale evitare la circolarità logica frequenzista: \textbf{la probabilità NON è il limite delle frequenze}. Se lo fosse, la Legge dei Grandi Numeri sarebbe una tautologia (una definizione travestita da teorema).
Invece, nell'assiomatica di Kolmogorov, la probabilità è una misura primitiva assegnata \emph{a priori}. La Legge dei Grandi Numeri ci dice che, sotto queste regole, il modello è internamente coerente: i risultati pratici a lungo termine si allineano con la costruzione matematica.
\end{hint}

\section{Esercizi Selezionati (Capitolo 2)}

Di seguito alcuni esercizi essenziali per fissare i concetti su variabili discrete, condizionamento e modelli classici.

\subsection*{Esercizio 1: Tempi d'attesa e parità}
Lanciando una moneta equa, qual è la probabilità che il primo successo avvenga in un lancio dispari?
\begin{spiegazione}
Sia $T \sim \text{Geom}(1/2)$. La probabilità è data dalla serie:
\[ \PP(T \in \{1,3,5,\dots\}) = \sum_{k=0}^\infty \PP(T = 2k+1) = \sum_{k=0}^\infty (1/2)^{2k+1} \]
Questo è pari a $\frac{1}{2} \sum (1/4)^k = \frac{1/2}{1 - 1/4} = \frac{1/2}{3/4} = \frac{2}{3}$.
\end{spiegazione}

\subsection*{Esercizio 2: Estrazioni con stop al primo disaccordo}
Due giocatori (A e B) estraggono (con rimpiazzo) da urne identiche ($r$ rosse, $b$ nere). Il gioco si ferma al primo turno in cui i risultati divergono. Chi ha estratto rosso, vince 1 Euro. Qual è la legge della vincita $V$ di A, e del numero di turni $N$? Sono indipendenti?
\begin{spiegazione}
La probabilità di successo al turno $n$ per A è $p = r/(r+b)$. Il gioco si ferma se (A estrae rosso e B nera) o viceversa. La probabilità che fermi a un turno è $2p(1-p)$.
$N$ ha legge geometrica di parametro $2p(1-p)$.
Dato che il gioco si è fermato, la probabilità che a vincere sia stato A (rosso) o B (rosso) è per simmetria esattamente $1/2$.
Quindi $V \sim B(1/2)$ e il risultato finale è \textbf{indipendente} dalla durata $N$.
\end{spiegazione}

\subsection*{Esercizio 3: Tempi di successo successivi}
In uno schema di Bernoulli con probabilità $p$, se so che il secondo successo è avvenuto al tempo $N$ ($T_2 = N$), qual è la legge del primo successo $T_1$?
\begin{spiegazione}
Usiamo Bayes: $\PP(T_1 = k | T_2 = N) \propto \PP(T_1 = k, T_2 = N)$.
L'evento $\{T_1 = k, T_2 = N\}$ equivale a: "1 successo al tempo $k$, insuccessi ovunque altro prima di $N$, e successo al tempo $N$". La probabilità è $p^2 (1-p)^{N-2}$, che \textbf{non dipende da $k$}.
Quindi, dato $T_2=N$, il tempo $T_1$ è uniformemente distribuito su $\{1, \dots, N-1\}$.
\end{spiegazione}

\subsection*{Esercizio 4: Estrazioni senza rimpiazzo (Variabili ipergeometriche)}
Estraggo $k$ palline su $n$ totali senza rimpiazzo. Qual è la probabilità di estrarre un elemento specifico $x$? E due elementi $x, y$?
\begin{spiegazione}
Sia $X_i$ l'esito della $i$-esima estrazione. Per simmetria, ogni estrazione marginalmente è uniforme: $\PP(X_i = x) = 1/n$.
La probabilità che $x$ sia estratto è $\PP(\cup_{i=1}^k \{X_i=x\}) = \sum \PP(X_i=x) = k/n$.
Per due elementi, la probabilità che vengano estratti entrambi è $k(k-1)/(n(n-1))$.
\end{spiegazione}

\subsection*{Esercizio 5: Somma di Bernoulli con numero di termini di Poisson}
Sia $N \sim \text{Poisson}(\lambda)$ indipendente da $X_i \sim B(p)$. Qual è la legge della somma parziale aleatoria $S_N = \sum_{i=1}^N X_i$? E la legge di $N - S_N$?
\begin{spiegazione}
Condizionando rispetto a $N$, otteniamo $S_N|N \sim \text{Bin}(N, p)$. Integrando, la mgf o le probabilità rivelano che \textbf{$S_N \sim \text{Poisson}(\lambda p)$} e \textbf{$N - S_N \sim \text{Poisson}(\lambda(1-p))$}.
Inoltre, sorprendentemente, $S_N$ e $N-S_N$ (i successi e i fallimenti totali) sono variabili \textbf{indipendenti}! Questo si verifica calcolando la mgf congiunta.
\end{spiegazione}

\subsection*{Esercizio 6: Strategia del raddoppio al gioco della moneta}
Se perdo raddoppio la puntata, se vinco mi fermo. Si inizia puntando 1. Se il capitale iniziale è 1023, mi rovino se perdo 10 volte di fila ($2^{10}-1 = 1023$).
Il guadagno netto, se si vince entro i 10 turni, è sempre esattamente $+1$. Se si perde, è $-1023$.
\begin{spiegazione}
La durata del gioco $T$ è ristretta a $\{1, \dots, 10\}$. Si ha $\PP(T=n) = q^{n-1}p$ per $n \le 9$, e $\PP(T=10) = q^9$.
Il valore atteso di $T$ si calcola con somme finite di serie geometriche derivate.
Questo esercizio ci ricorda che le "strategie infallibili" come il raddoppio delle puntate collassano di fronte alla finitezza dei capitali, portando a rischi di rovina asimmetrici enormi.
\end{spiegazione}
